\subsection{Deleting both companion parities in a growing character range}
\label{sec:growing-component-window}

Write $T=M_{\Lambda\text{-tors}}$ for the torsion submodule of the
actual ramified module used in Theorem~\ref{thm:odd-torsion-cubic}.
Here the parity of a companion and the parity of its reflected
unramified class component must be distinguished.

\begin{proposition}[K17-BOTH-COMPANIONS]
\label{prop:both-companions}
Let $p\geq5$ and let $a$ be even with $2\leq a\leq p-3$.
If $p\nmid\operatorname{num}(B_a)$, then for every $n\geq0$
\[
 A_n^{[a]}=A_n^{[1-a]}=0,\qquad
 X^{[a]}=X^{[1-a]}=0,\qquad
 M^{[a]}=0,\qquad M^{[1-a]}\simeq\Lambda.
 \tag{BC.1}
\]
In particular $T^{[a]}=T^{[1-a]}=0$. Thus for every
$p>2^{31}$ both companion parities in (BC.1) vanish for every
even $a$ from two through thirty-two. This last assertion uses
the exact Bernoulli certificate (PC.3), independently of Alp\"oge.
\end{proposition}

\begin{proof}
The Herbrand implication gives $A_0^{[1-a]}=0$:
the exponent $p-a$ in the usual statement is congruent to
$1-a$ modulo $p-1$. Its reflection, in the exact form proved
in Lemma~\ref{lem:isolated-reflection}, gives $A_0^{[a]}=0$.
Both characters satisfy the all-character finite-layer implication
(CC.1), so both class components vanish at every layer and in
the norm inverse limit.

Let $A_\infty=\varinjlim_n A_n$, with extension-of-ideal maps.
Its component $[1-a]$ is zero because every term in that
component is zero. The exact Kummer-duality isomorphism
\[
                    (A_\infty^-)^\vee(1)\simeq M^+
\]
of \cite[Proposition~3.4.7 and its proof, p.~92]{SharifiIwasawaFull2026}
therefore gives $M^{[a]}=0$ integrally. Indeed duality negates
the character and the Tate twist adds one, so $[1-a]$ becomes
$[-(1-a)+1]=[a]$. This uses the direct limit explicitly;
it does not confuse it with the norm inverse limit $X$.

For the other parity, (CC.2) applied to $[a]$ gives
$M^{[1-a]}\simeq\Lambda$ unless $a=2$. The exceptional
component is then $M^{[-1]}$, whose integral rank-one
freeness is the Coleman calculation already used in
Corollary~\ref{cor:odd-companion-exclusion}.
The asserted torsion vanishings follow immediately.
Finally (PC.3) proves that all numerator prime divisors for
$2\leq a\leq32$ are smaller than $2^{31}$; for a prime above
this bound the hypotheses apply to every such $a$.
The even-companion conclusion is not extended to the small
numerator primes by Vandiver: an odd class component can be
nonzero at a Vandiver prime, as the established example at
$p=37$, $a=32$, illustrates.
\end{proof}

\begin{sourcefact}[Partial regularity for almost all primes]
\label{source:partial-regularity}
Fix $\alpha>1/2$ and set
\[
                  b_\alpha(p)=
                  \left\lfloor\frac{\sqrt p}{(\log p)^\alpha}\right\rfloor.
\]
There is a set $\mathcal P_\alpha$ of primes such that
$p\nmid\operatorname{num}(B_a)$ for every $p\in\mathcal P_\alpha$
and every even $2\leq a\leq\min(b_\alpha(p),p-3)$, and
\[
 \#\{p\leq x:p\notin\mathcal P_\alpha\}
          =O_\alpha\!\left(\frac{x}{(\log x)^{2\alpha}}\right).
 \tag{BC.2}
\]
This is \cite[Theorem~1.1, p.~2; full proof in \S3,
pp.~9--10, using \S2, pp.~6--8]{ChenEtAl2026}.
Its reflected base-class consequence is Corollary~1.4, p.~3.
The source's formalization is cited but has not been executed
in this run; the complete mathematical proof of the stated
counting theorem has been read.
\end{sourcefact}

\begin{corollary}[K17-GROWING-COMPONENT-WINDOW]
\label{cor:growing-component-window}
For every $p\geq5$ in $\mathcal P_\alpha$ and every even
$2\leq a\leq\min(b_\alpha(p),p-3)$, all conclusions (BC.1)
hold. In particular, neither the even companion $T^{[a]}$
nor the reflected odd companion $T^{[1-a]}$ occurs in the
actual arithmetic group. For every positive integer $k$ with
\[
                     k\equiv-a/2\pmod{(p-1)/2}
\]
one has $K_{4k}(\mathbf Z)\{p\}=0$, and the full-domain
Soul\'e character $\kappa_{2k+1,p}$ is onto $\mathbf Z_p$.
These conclusions hold on a density-one set of primes with
exception bound (BC.2), independently of Alp\"oge.
\end{corollary}

\begin{proof}
Apply Proposition~\ref{prop:both-companions} separately to each
index in the supplied simultaneous Bernoulli range. Its hypotheses
hold simultaneously by the definition of $\mathcal P_\alpha$,
so the exceptional-prime set does not grow when passing to all
layers or to both companion parities. Since $2\alpha>1$,
division of (BC.2) by $\pi(x)\sim x/\log x$ gives relative
prime density zero for the exceptions.
For the displayed congruence one has $-2k\equiv a\pmod{p-1}$.
The exact comparison (K2) therefore gives the primary $K$-group
vanishing from $A_0^{[a]}=0$. The full-domain index formula
(K18), with $m=2k+1$, gives the Soul\'e conclusion.
\end{proof}

The growing range removes specified actual generators from the
remaining LM6 A.45 gate. It does not exhaust the characters of
$\Delta$ or assert Greenberg, Vandiver, or the entire kernel
equality on this density-one set. For the universal fixed range,
small primes retain the already proved full-kernel treatment;
the new assertion about absent even companions is reserved for
the explicitly stated large primes.
